% ============================================================
\section{Ejercicio 3}

% ============================================================

\begin{tcolorbox}[enunciado]
Prueba que la colección de lenguajes decidibles en cerrada bajo las siguientes operaciones:
\begin{enumerate}[label=(\alph*)]
    \item Unión.
    \item Intersección.
    \item  Concatenación.
    \item  Complemento.
\end{enumerate}
¿Para cuáles de las operaciones anteriores se mantiene la cerradura para las clases $P,NP$?
\end{tcolorbox}

% ------------------------- Solución -------------------------
\subsection{Respuestas}

\subsection*{a}

Para demostra que $L_3(M_3)=L_1 \cup L_2$, suponemos lo siguiente.\\

Suponemos que $L_1$ es decidible si exite una maquina $M_1$ tal que $L(M_1) = L_1$ y $M_1$ es total. De la misma forma suponemos un $L_2$ que es decidible si exite una maquina $M_2$ tal que $L(M_2) = L_2$ y $M_2$ es total.\\

Por lo tanto lo que se quiere demostrar es que existe una $M_3$ tal que $L_3(M_3) = L_1 \cup L_2$ y $M_3$ es total.\\

Suponemos el funcionamiento de $M_3 = U$ como el siguiente:
\begin{enumerate}
    \item U simula $M_1$ con x
        Si $M_1(x) = 1$, entonces U acepta
    \item U simula $M_2$ con x
        Si $M_2(x) = 1$, entonces U acepta
    \item U en otro caso rechaza
\end{enumerate}

Si suponemos que $M_1(x)=1$, entonces $L_1(M_1)$ acepta, dada la union, si $x$ bajo la ejecuion de $M_1$ acepta, entonces $L_1 \cup L_2$ se cumple y acepta.\\

Si suponemos que $M_2(x)=1$ acepta, entonces $L_2(M_2)$ acepta, dada la union, si $x$ bajo la ejecuion de $M_2$ acepta, entonces $L_2 \cup L_1$ se cumple y acepta.\\

Si $M_1(x) = 0$, entonces $L_1(M_1)$ rechaza, dada la union, $M_2(x)=1$, por lo tanto $L_1 \cup L_2$ se cumple. Sucede lo mismo para $M_2$\\

Por doble contencion tenemos que $L_3 = L_1 \cup L_2$.\\

Para analizar la complejidad temporal de $M_3$ con una entrada $x$ de longitud $n = \vert{}x\vert{}$:
Sea $t_1(n)$ el tiempo que tarda $M_1$ en procesar una entrada de longitud $n$.\\
Sea $t_2(n)$ el tiempo que tarda $M_2$ en procesar una entrada de longitud $n$.\\

Como $M_3$ ejecuta $M_1$ y (en el peor de los casos) posteriormente $M_2$ en secuencia sobre la misma entrada $x$, el tiempo de ejecución total $T_{M_3}(n)$ está acotado por la suma de los tiempos de ambas máquinas:
$$T_{M_3}(n) = O(t_1(n) + t_2(n))$$

Caso Polinomial:
Si $M_1 \in TIME(n^{k_1})$ y $M_2 \in TIME(n^{k_2})$ para constantes $k_1, k_2 \ge 1$:$$T_{M_3}(n) = O(n^{k_1} + n^{k_2}) = O(n^{\max(k_1, k_2)})$$

Pero como solo se quiere saber si es decidible o no, no es necesario analizar su complejidad temporal, pues solo debemos garantizar en aceptar o rechazar en un número finito de pasos para toda entrada $x$.

\subsection*{b}

Para demostra que $L_3(M_3)=L_1 \cap L_2$, suponemos lo siguiente.\\

Suponemos que $L_1$ es decidible si exite una maquina $M_1$ tal que $L(M_1) = L_1$ y $M_1$ es total. De la misma forma suponemos un $L_2$ que es decidible si exite una maquina $M_2$ tal que $L(M_2) = L_2$ y $M_2$ es total.\\

Por lo tanto lo que se quiere demostrar es que existe una $M_3$ tal que $L_3(M_3) = L_1 \cap L_2$ y $M_3$ es total.\\

Suponemos el funcionamiento de $M_3 = U$ como el siguiente:
\begin{enumerate}
\item $M_3$ simula a $M_1$ con la entrada $x$
\item Si $M_1(x) = 0$ rechaza, $M_3$ rechaza 
\item Si $M_1(x) = 1$ acepta, entonces $M_3$ simula a $M_2$ con la entrada $x$
\item Si $M_2(x) = 1$ acepta, $M_3$ acepta
\item Si $M_2(x) = 0$ rechaza, $M_3$ rechaza 
\end{enumerate}

Si suponemos que $M_1(x)=1$, entonces $L_1(M_1)$ es reconocible,de esta manera para la misma $x$,$M_2(x)=1$ acepta; $L_1(M_1) \cap L_2(M_2)$ se cumple.\\

Si suponemos que $M_2(x)=1$, entonces $L_2(M_2)$ es reconocible, entonces para la misma $x$,$M_1(x)=1$ acepta, de igual manera es reconocible; $L_1(M_1) \cap L_2(M_2)$ se cumple.\\

Dado que $M_1$ es una máquina total, la primera simulación siempre finaliza en un número finito de pasos. Si $M_1$ rechaza, $M_3$ termina inmediatamente. Si $M_1$ acepta, se pasa a simular $M_2$, la cual también es total y garantiza terminar en un número finito de pasos. En ningún caso $M_3$ entra en un bucle infinito, demostrando que $M_3$ es total.\\

Por lo tanto, $L(M_3) = L_1 \cap L_2$ y $M_3$ es total, lo que demuestra que la clase de lenguajes decidibles es cerrada bajo intersección.\\

\subsection*{c} 
Para demostra que $L_3(M_3)=L_1 L_2$, suponemos lo siguiente.\\

Suponemos que $L_1$ es decidible si existe una maquina $M_1$ tal que $L(M_1) = L_1$ y $M_1$ es total. De la misma forma suponemos un $L_2$ que es decidible si existe una maquina $M_2$ tal que $L(M_2) = L_2$ y $M_2$ es total.\\

Por lo tanto lo que se quiere demostrar es que existe una $M_3$ tal que $L_3(M_3) = L_1 L_2$ y $M_3$ es total. Para esto suponemos una maquina de turing tal que $M_3$ simule a $M_1$ y $M_2$, así. Suponemos que $L_3= \{x=yz: y \in L_1 , z \in L_2 \}$

\begin{enumerate}
\item Para cada posible división de la cadena $x$ en dos partes $y$ y $z$ tales que $x = yz$ (desde $y = \epsilon, z = x$ hasta $y = x, z = \epsilon$):
\begin{enumerate}
\item Simular $M_1$ con la entrada $y$
\item Simular $M_2$ con la entrada $z$
\item Si $M_1(y) = 1$ y $M_2(z) = 1$, entonces $U$
\end{enumerate}
\item Si $U$ prueba todas las posibles divisiones $x = yz$ y en ninguna se cumple para $M_1$ y $M_2$ ambas acepten, entonces $M_3$ rechaza.
\end{enumerate}

Dada la cadena de entrada, se simula la primera parte de $x$ en $M_1$ hasta alcanzar un estado final y, a continuación, se pasa a $M_2$ para completar la segunda mitad del cómputo.\\

Puesto que una cadena $x$ de longitud $n$ solo tiene un número finito de particiones las cuales son $n+1$ posibles divisiones y $M_1$ y $M_2$ son máquinas totales que siempre se detienen, el bucle de $M_3$ siempre terminará en un número finito de pasos. Entonces la concatenacion de los lenguajes con sus respectivas maquinas como $L_1(M_1)$ y $L_2(M_2)$;  $L_1L_2$ es cerrado bajo la operacion de concatenacion. 

\subsection*{d}

Para demostrar que $L_3(M_3)= (L_1(M_1))^c$\\

Por definición, sabemos que una maquina es deidible si existe una $M_3$ tal que $L_3(M_3) = L_1(M_1)^c$. Asi, suponemos el siguiente funcionamiento de la maquina $M_3 = U$.\\

\begin{enumerate}
    \item U simula $M_1$ con x. Si $M_1(x) = 0$, entonces $U$ acepta
    \item U simula $M_1$ con x. Si $M_1(x) = 1$, entonces $U$ rechaza
    \item $U$ rechaza en otro caso
\end{enumerate}

Dado que $M_1$ es una máquina total, para cualquier entrada $x \in \Sigma^*$ solo existen dos posibilidades:\\

Si $x \in (L(M_1))^c$, entonces $x \notin L(M_1)$, lo que significa que $M_1(x) = 0$. Por el paso (1), $M_3$ acepta $x$. Si $x \notin (L(M_1))^c$, entonces $x \in L(M_1)$, lo que significa que $M_1(x) = 1$. Por el paso (2), $M_3$ rechaza $x$.\\

Dado que $M_3$ siempre se detiene e invierte exactamente las respuestas de $M_1$, se cumple que $L(M_3) = (L(M_1))^c$ y $M_3$ total. Por lo tanto, el complemento de cualquier lenguaje decidible es también decidible.

\subsection*{e}¿Para cuáles de las operaciones anteriores se mantiene la cerradura para las clases P, NP?\\
